\documentclass[11pt]{article}

\usepackage[margin=1in]{geometry}
\usepackage{amsmath}
\usepackage{booktabs}
\usepackage{array}
\usepackage{graphicx}
\usepackage{hyperref}

\title{Comparison Tools for Skipper-CCD Traces:\\Parametric Fits and Summary Measurements}
\author{ML-assisted optimization of CCD operating parameters --- working note}
\date{October 2026}

\newcommand{\code}[1]{\texttt{#1}}

\begin{document}
\maketitle

\begin{abstract}
This note describes all 39 tools used by \code{compare\_trace.py} to decide whether two raw
video-channel traces of a skipper CCD can be told apart: 11 parametric fits (time constants and amplitudes
of small physical models), 19 summary measurements of the skip region, and 9 measurements of the horizontal
and vertical transfer blocks. Of the 39, 35 enter the verdict; the 4 vertical-block tools are printed for
information only. The note explains what each tool does and how it is computed, how the numbers are compared
between two traces, and what the tests on real data showed: channel~0 against channel~3, and repeats of channel~0 with identical
parameters (Sections~\ref{sec:results} and~\ref{sec:repeats}). All numbers are in
uncalibrated ADU unless stated otherwise. The tools have so far been tried only on traces with
\emph{identical} operating parameters; they have not yet been shown to respond to parameter changes.
\end{abstract}

\section{Purpose and context}

The ML program needs tools that distinguish traces taken with different operating parameters. Simple
summary numbers (edge heights, window means, noise power; Section~\ref{sec:other}) were implemented first. The fit
tools (Section~\ref{sec:fits}) add a second kind of measurement: instead of reading a value off the waveform, they fit a small physical
model and report its parameters. Time constants in particular are expected to depend on bandwidth and
timing rather than on the gain of an individual amplifier, which makes them candidates for tools that
are less sensitive to amplifier-to-amplifier differences. That expectation is a hypothesis; it has not
been tested against parameter changes.

\section{Data preparation shared by all fits}

A trace is decoded (\code{decode\_trace.py}) and split into regions (\code{trace\_regions.py}). The fits
described here work on the \emph{skip region} of one pixel. Let $X_{kj}$ be the raw ADC value of sample
$j$ in skip cycle $k$. For the traces analysed so far, a cycle has $P=489$ samples (32.6~$\mu$s at
15~MHz), a pedestal window of $n_{\rm ped}=150$ samples at the start of the cycle, and a signal window
starting at sample $j_{\rm sig}=207$.

\begin{description}
  \item[Baseline.] $b_k = \frac{1}{n_{\rm ped}}\sum_{j<n_{\rm ped}} X_{kj}$, the mean of the pedestal
        window of skip $k$. It walks slowly across the skips (reset-leakage droop).
  \item[Baseline-removed cycle.] $s_{kj} = X_{kj} - b_k$.
  \item[Template.] $m_j = \frac{1}{K-5}\sum_{k\ge 5} s_{kj}$, the mean shape of the cycle. The first five
        skips are excluded because they behave differently.
  \item[CDS.] $c_k$ = pedestal mean $-$ signal mean of skip $k$, taken from the inline level-1 records
        written by the firmware. $c_{\ge 10}$ denotes the CDS values of skips $10,\dots,K-1$.
\end{description}

Each of the 10 pixels of a trace gives one set of fit results; the pixels are the replicates used in the
statistical comparison (Section~\ref{sec:stats}).

\section{The fit tools}
\label{sec:fits}

\subsection{Exponential fit by variable projection}
\label{sec:exp}

Three tools (settling after the pulse, dip recovery, baseline droop) share one routine,
\code{fit\_exponential}, which fits
\begin{equation}
  y(x) = c + A\,e^{-x/\tau}.
\end{equation}
The parameters $A$ and $c$ enter linearly, so for any trial $\tau$ they are found exactly by linear
least squares with basis $\{e^{-x/\tau},\,1\}$. Only $\tau$ is searched:
\begin{enumerate}
  \item 60 trial values, log-spaced between $\tau_{\min}$ and $\tau_{\max}$; keep the one with the
        smallest sum of squared residuals.
  \item 41 values, log-spaced over $[\tau^*/1.25,\;1.25\,\tau^*]$ around that point.
  \item A parabola is fitted through the minimum and its two neighbours in $(\ln\tau,\ \mathrm{SSE})$,
        and its vertex is taken as $\tau$. This removes the quantization that the grid would otherwise
        impose on the reported time constants. (If the minimum sits at the edge of the fine grid, the grid
        value is kept.)
  \item $A$ and $c$ are recomputed by least squares at the final $\tau$.
\end{enumerate}
No gradient-based optimizer is used, so the fit cannot fail to converge; it can, however, return a value at
the edge of the search range if the data do not contain an exponential. The routine returns NaN if fewer
than five points are given or the data are not finite; such replicates are dropped from the test.

\subsubsection{Settling after the pulse (\code{fit: settling time after pulse}, \code{fit: settling amplitude after pulse})}

\textbf{What it measures.} Every skip cycle contains a large positive pulse after the signal window
(visible in \code{plots/regions\_templates.png}: the pulse peaks at about 25~$\mu$s). After the peak the
video signal relaxes back to its baseline. The relaxation is expected to reflect the finite bandwidth
of the source follower and preamplifier, so its time constant should depend on the amplifier chain and
on timing, and its amplitude on the size of the clock feedthrough.

\textbf{How.} The template $m_j$ is cut at its maximum after the signal window,
$j_{\rm peak} = j_{\rm sig}+n_{\rm ped}+\arg\max_j m_j$, and $m_j$ for $j\ge j_{\rm peak}$ is fitted against
$x=(j-j_{\rm peak})/15$ in $\mu$s with $\tau\in[0.05,\,20]\ \mu$s. Reported: $\tau$ and $A$.

\subsubsection{Dip recovery (\code{fit: dip recovery time}, \code{fit: dip recovery amplitude})}

\textbf{What it measures.} Between the pedestal and the signal window the template shows a short
negative dip (the summing-well/output-gate edge) followed by recovery. Its amplitude measures the size of
that feedthrough; its time constant the speed of the recovery.

\textbf{How.} $j_{\rm dip} = n_{\rm ped}+\arg\min_j m_j$ over $j\in[n_{\rm ped},\,j_{\rm sig})$. The template from the
minimum up to the start of the signal window is fitted with $\tau\in[0.02,\,10]\ \mu$s. Reported: $\tau$ and
$A$ (negative, since the dip points downward).

\subsubsection{Baseline droop (\code{fit: baseline droop time constant}, \code{fit: baseline droop amplitude})}

\textbf{What it measures.} The pedestal baseline $b_k$ rises across the skips of a pixel (in the sample trace
from about $-22\,000$ to $-3\,500$ ADU over 500 skips). The node voltage creeps toward a limiting value; a
saturating exponential describes it well. The time constant (in skips) is the characteristic time of this
creep, which is related to the leakage current and the node capacitance (Eq.~(24) of the noise-landscape
note); the amplitude is the total excursion.

\textbf{How.} $b_k$ for $k=1,\dots,K-1$ (skip 0 is excluded) is fitted against $x=k$ with
$\tau\in[5,\,5000]$ skips. Reported: $\tau$ (skips) and $A$ (ADU, negative for a rising baseline).

\subsection{Damped ringing of the first skips (\code{fit: first-skip ringing ...})}

\textbf{What it measures.} The CDS of the first skips of a pixel does not settle immediately: in the sample
trace skip~0 is about 320~ADU above the later values, followed by swings of changing sign
(\code{plots/skip\_evolution.png}, panel C). This tool describes that transient with a damped oscillation.

\textbf{Model (corrected 2026-10-13).} For $y_k = c_k - \mathrm{median}(c_{\ge 10})$, $k=1,\dots,19$:
\begin{equation}
  y_k = C + e^{-k/\tau}\bigl(a\cos 2\pi f k + b\sin 2\pi f k\bigr).
\end{equation}
Skip~0 is left out (its larger transient is the separate tool ``first-skip CDS excess'') and the constant $C$ absorbs the
level that has not settled within 20 skips. $f$ is searched in $0.02$--$0.5$ cycles/skip and $\tau$ in $0.5$--$30$ skips,
then refined; $C$, $a$, $b$ are solved by least squares. On the 2026-10-08 and 2026-10-12 traces this gives
$f \approx 0.23$ cycles/skip (a period of about 4.3 skips), $\tau \approx 1.5$--$3$ skips, in every pixel.

\textbf{The original model} (used for the numbers in this note) had no $C$ and included skip~0,
$y_k = e^{-k/\tau}\bigl(a\cos 2\pi f k + b\sin 2\pi f k\bigr)$ for $k = 0,\dots,19$. Its lowest non-zero frequency
imitated a decay with an offset, so it returned $f = 0.01$ for almost every pixel. The frequency $f$ (51 values from 0 to 0.5 cycles per skip; $f=0$ is a pure decay) and the decay length
$\tau$ (25 log-spaced values from 1 to 30 skips) are searched on a grid; $a$ and $b$ are linear and solved by
least squares for each pair. The pair with the smallest squared residual is kept.

\textbf{Reported.} Frequency $f$ (cycles per skip), decay length $\tau$ (skips), and amplitude
$\sqrt{a^2+b^2}$ (ADU).

\textbf{Caveat.} Twenty points and four effective parameters are few. In the sample data the fitted frequency
differed strongly between channels (0.03 against 0.17 cycles per skip), and for several pixels the best fit
is essentially a pure decay. The decay length and amplitude are more trustworthy than the frequency.

\subsection{Noise model (\code{fit: white noise per skip}, \code{fit: correlated noise floor})}

\textbf{What it measures.} Averaging $N$ skips reduces uncorrelated (white) noise like $1/\sqrt{N}$ but
cannot reduce correlated noise (1/$f$, drifts). The noise of an $N$-skip average is therefore expected to
follow
\begin{equation}
  \sigma(N)^2 = \frac{\sigma_{\rm w}^2}{N} + \sigma_{\rm f}^2 ,
\end{equation}
where $\sigma_{\rm w}$ is the white noise per skip and $\sigma_{\rm f}$ the floor that averaging cannot remove.

\textbf{How.} The CDS values $c_{\ge10}$ are detrended by subtracting a degree-4 polynomial in the skip number (this removes
the slow drift and keeps the scatter; a short running mean must not be used, because it also cancels the average over
blocks of the same length). For $N\in\{1,2,5,10,25,50\}$ the detrended values are divided into
non-overlapping blocks of $N$ skips, and $\sigma(N)$ is the sample standard deviation (ddof $=1$) of the block
means. $\sigma_{\rm w}^2$ and $\sigma_{\rm f}^2$ follow from a linear least-squares fit of $\sigma(N)^2$ against
$(1/N,\,1)$; negative coefficients are set to zero. Reported: $\sigma_{\rm w}$ and $\sigma_{\rm f}$ in ADU.

\textbf{Caveat.} The model assumes the skip-to-skip noise is white plus a constant floor. The measured noise of
$N$-skip averages in the sample traces falls \emph{faster} than $1/\sqrt{N}$ (for example 0.34~ADU at $N=50$ against
0.81~ADU expected from $\sigma(1)=5.7$~ADU), because the CDS values of nearby skips are anti-correlated (lag-2
autocorrelation about $-0.2$). A floor cannot describe that, so the fitted floor is small (0.16 and 0.22~ADU on the two
channels) and the model fits the data poorly; $\sigma_{\rm w}$ is close to $\sigma(1)$. The lag-2 autocorrelation
is therefore also measured directly (Section~\ref{sec:other}).

\section{The other tools (summary measurements)}
\label{sec:other}

The remaining 28 tools are direct measurements on the template, the baseline and the CDS values. They use
the quantities defined in Section~2, with $n_{\rm ped}=150$ the pedestal-window length, $j_{\rm sig}$ the start of the
signal window and $P$ the cycle length. Curves are block-averaged to at most 500 points before testing.

\subsection{Skip region (19 tools, one value or curve per pixel)}

\begin{description}
  \item[mean shape] The template $m_j$ as a curve.
  \item[baseline droop curve] $b_k-b_1$ for $k\ge1$, as a curve over the skips.
  \item[dip height] $\min_j m_j$ over $j\in[n_{\rm ped},\,j_{\rm sig})$.
  \item[signal-window level] mean of $m_j$ over the signal window, excluding 10 samples at each end.
  \item[pulse peak height] $\max_j m_j$ over $j\ge j_{\rm sig}+n_{\rm ped}$.
  \item[end-of-cycle tail level] mean of the last 30 samples of $m_j$.
  \item[baseline at skip 1] $b_1$.
  \item[droop total] $b_{K-1}-b_1$.
  \item[droop half-time] the number of skips (counted from skip 1) after which the baseline has covered half of
        its total change.
  \item[CDS mean] mean of $c_{\ge10}$.
  \item[CDS rms (detrended)] sample standard deviation of $c_{\ge10}$ after subtracting a degree-4 polynomial in the skip number.
  \item[first-skip CDS excess] $c_0-\mathrm{median}(c_{\ge10})$.
  \item[skips 1--4 CDS ringing] root-mean-square of $c_k-\mathrm{median}(c_{\ge10})$ for $k=1,\dots,4$.
  \item[CDS autocorrelation at lag 2] correlation between $c_k$ and $c_{k+2}$ of the detrended CDS values
        ($c_{\ge10}$): zero for white noise, about $-0.2$ in the sample traces.
  \item[noise of 10-skip and 50-skip average] $\sigma(N)$ for $N=10$ and $N=50$, as defined for the noise model in
        Section~\ref{sec:fits}.
  \item[noise power in three bands] For skips $k\ge5$ the pedestal-window samples of $s_{kj}-m_j$ are taken, the
        mean of each window is removed, and the squared magnitude of its Fourier transform is averaged over
        the skips and divided by $n_{\rm ped}$. The tool is $\log_{10}$ of the summed power in
        $0.1$--$1$, $1$--$3$ and $3$--$7.5$~MHz. The frequency resolution is 100~kHz.
\end{description}

\subsection{Horizontal block (5 tools, one value per pixel)}

For each pixel, all (typically 19) horizontal transfers are stacked, the mean of the first 30 samples of each
transfer is subtracted, and $h_j$ denotes the mean over the transfers.

\begin{description}
  \item[mean shape] $h_j$ as a curve.
  \item[minimum, maximum] $\min_j h_j$ and $\max_j h_j$.
  \item[end level] mean of the last 50 samples of $h_j$.
  \item[scatter between transfers] median over time of the standard deviation across transfers.
\end{description}

\subsection{Vertical block (4 tools, one value per transfer; information only)}

Same definitions as the horizontal block (mean shape, minimum, maximum, end level), but one replicate is one
vertical transfer (typically 18 per trace) and no scatter tool exists. These tools are excluded from the verdict
because consecutive transfers drift along the block and are not interchangeable (Section~\ref{sec:stats}).

\subsection{What the tools measure: an annotated example}

Figure~\ref{fig:annotated} shows one pixel of a healthy trace with the quantity each tool reports marked on the data
(panels A, D, E, G, H), the fits drawn over the data (B, C, D, E, F), and the replicate structure used in the statistics (I).
Two things can be read off directly. First, the single-exponential fits in panels B and C describe the dip recovery and
the settling only approximately: the data have a slower tail than one exponential, so the fitted time constants are
descriptive parameters, not circuit values. Second, panel F shows the noise of $N$-skip averages falling faster than the
white-plus-floor model, as discussed in the noise-model caveat.

\begin{figure}[htbp]
  \centering
  \includegraphics[width=\textwidth]{../plots/tools_annotated.png}
  \caption{What the tools measure (ACM 101 channel 0, pixel 6). A: skip-cycle template with the pedestal and signal
  windows, the dip height, the pulse peak height, the signal-window level and the end-of-cycle tail level.
  B, C: exponential fits to the recovery after the dip and to the settling after the pulse. D: baseline of every skip with
  the droop fit, the baseline at skip 1, the total droop and the half-time. E: CDS of the first skips, the first-skip excess,
  the damped-ringing fit and the CDS mean. F: noise of $N$-skip averages with the white-plus-floor fit and the $1/\sqrt{N}$
  line. G: noise power of the pedestal window with the three frequency bands. H: horizontal transfer with minimum,
  maximum and end level. I: the pulse peak height of each usable pixel (dots) and of each run (bars) for four repeat runs.}
  \label{fig:annotated}
\end{figure}

\section{How two traces are compared}
\label{sec:stats}

Each tool gives one number per pixel (the \emph{replicate}). For traces A and B:
\begin{itemize}
  \item \textbf{Statistic.} Welch's $z=(\bar x_A-\bar x_B)/\sqrt{s_A^2/n_A+s_B^2/n_B}$; for a curve, the
        root-mean-square of $z$ over the curve.
  \item \textbf{Permutation test.} The A/B labels of the replicates are shuffled; the $p$-value is the
        fraction of labelings that give a statistic at least as large as the real one. If there are at most 50\,000 distinct
        labelings they are all enumerated and $p$ is exact (8 usable pixels per trace give 12\,870 labelings, so the smallest possible
        $p$ is $2/12\,870=1.6\times10^{-4}$: the real labeling and its mirror image); otherwise 10\,000 random labelings are used and
        the smallest reportable value is $10^{-4}$. Many strongly different tools therefore tie at this floor and can only be
        ranked by $z$.
  \item \textbf{Threshold.} A tool \emph{distinguishes} the traces if $p<\alpha/N_{\rm tools}$ with
        $\alpha=0.01$ (Bonferroni). With 35 tools this is $2.9\times10^{-4}$. Tools with $p<0.05$ but above
        the threshold are called \emph{marginal}.
  \item \textbf{Controls.} The same test is run on trace A alone, split into even against odd pixels and into
        first against second half (row~1 against row~2). If any tool passes the corrected threshold in a
        control, the verdict is declared inconclusive.
  \item \textbf{Vertical block.} The 18 vertical transfers are not interchangeable (a control showed their first
        and second halves differ), so the vertical tools are printed for information and excluded from the
        verdict.
\end{itemize}
Because the fitted time constants scatter very little between pixels (the template is the mean of 500 skips),
$z$ becomes very large for them even for small differences. Rank tools by $p$; treat $z$ as a rough size only.

\section{First results: ACM 101, channel 0 against channel 3}
\label{sec:results}

Both traces use identical operating parameters; they differ only in which amplifier was read. The verdict was
that the traces \emph{are} distinguishable, and both controls passed. The fits, averaged over the 10 pixels:

\begin{center}
\begin{tabular}{l r r l}
\toprule
Fit & ch 0 & ch 3 & $p$ (verdict) \\
\midrule
Settling time after pulse ($\mu$s) & 0.622 & 0.578 & $3.0\times10^{-4}$ (marginal) \\
Settling amplitude after pulse (ADU) & 43\,090 & 36\,420 & $<10^{-4}$ (distinguishes) \\
Dip recovery time ($\mu$s) & 0.316 & 0.398 & $2.0\times10^{-4}$ (distinguishes) \\
Dip recovery amplitude (ADU) & $-$12\,140 & $-$9\,549 & $<10^{-4}$ (distinguishes) \\
Baseline droop time constant (skips) & 237 & 241 & $2.0\times10^{-4}$ (distinguishes) \\
Baseline droop amplitude (ADU) & $-$19\,420 & $-$17\,970 & 0.35 (no) \\
First-skip ringing frequency (cyc/skip) & 0.03 & 0.173 & $4.5\times10^{-3}$ (marginal) \\
First-skip ringing decay (skips) & 1.25 & 1.77 & $7.0\times10^{-4}$ (marginal) \\
First-skip ringing amplitude (ADU) & 5\,517 & 1\,464 & $1.8\times10^{-3}$ (marginal) \\
White noise per skip (ADU) & 5.75 & 6.89 & $<10^{-4}$ (distinguishes) \\
Correlated noise floor (ADU) & 0.163 & 0.222 & 0.76 (no) \\
\bottomrule
\end{tabular}
\end{center}

\noindent Values are means over the 10 pixels. The ringing amplitude is pulled up by a few pixels where the
fit is poorly constrained, so its mean is less representative than the other rows. All $p$-values are in the
CSV written by the script (\code{plots/compare\_ch0\_vs\_ch3\_acm101.csv}).

\paragraph{Observations.} The two amplitude fits (dip and settling) were among the most strongly separating tools
overall, as expected if the amplitudes scale with the amplifier gain and coupling. The time constants differ much less in
relative terms: about $2\%$ for the droop, $7\%$ for the settling and $26\%$ for the dip recovery. Smaller is not the same
as absent: Section~\ref{sec:repeats} shows that the settling and dip-recovery times are so reproducible between repeat
runs that even their few-percent differences between the two amplifiers are far outside the repeat noise. Whether the time
constants respond to operating parameters, and by how much more than to the amplifier, is still to be tested with
parameter-scan data.

\section{Repeats of one amplifier against two different amplifiers}
\label{sec:repeats}

A session on 2026-10-08 provides five traces of channel~0 and five of channel~3 of module 101, all taken with identical
operating parameters (four of the channel-0 traces within 12 seconds, the fifth 90 seconds later). This is the first real
test of whether the tools say ``same'' for repeats.

\paragraph{Data quality.} One channel-0 trace lost data during the acquisition: it has about 2\% fewer words than the others, 183 fewer
inline records, pedestal windows of unequal length, and its first pixel has only 412 skips. The keyword framing is intact, so
this looks like dropped network packets, not a decoding problem. \code{DecodedTrace.quality\_problems()} now detects it, and
\code{compare\_trace.py} refuses damaged traces unless told otherwise. The figures below use the four healthy channel-0
traces and the five channel-3 traces.

\paragraph{The pixel-level test is overconfident for repeats.} Applying the test of Section~\ref{sec:stats} to pairs of
channel-0 repeats flagged 0 to 8 of the 35 tools as different (per pair: 6, 8, 0, 6, 6, 8) even though nothing was changed, against 14 to 19 for
pairs of different amplifiers (left panel of Figure~\ref{fig:verdicts}). The causes found in the data:
\begin{itemize}
  \item The first pixel of one run held a very large signal (CDS mean 18\,680~ADU against 20--30~ADU normally) and its pulse peak and dip were distorted.
  \item The first pixel of each row (pixels 0 and 5) always differs: its baseline at skip 1 is about $-33\,000$~ADU against about $-19\,000$~ADU.
  \item The CDS level changed between runs with identical parameters (about 25, 80, 290, 25~ADU in time order; most likely real charge, not
        confirmed), which makes the CDS mean and the signal-window level useless as parameter tools.
  \item All pixels in a run share a small run-wide shift (about 0.16\% of the pulse peak, 74~ADU), larger than the 9~ADU error derived from
        the pixel scatter, because pixels within a run do not vary independently.
\end{itemize}
\code{usable\_pixels} in \code{compare\_trace.py} excludes the row-start pixels and charge outliers (median/MAD rule) and is available as
\code{--clean-pixels}; the figures use it.

\paragraph{The right noise floor is the run-to-run spread.} For every tool the spread of the per-run value over the channel-0 repeats was
computed. The right panel of Figure~\ref{fig:z} shows the channel-0 to channel-3 difference in units of that spread. The largest are the
dip-recovery time (358), the settling time (160), the pulse peak height (71), the dip-recovery amplitude (71), the dip height (63) and the
noise power at 1--3~MHz (59). Seven tools are below the factor of 3 that I take as ``detectable'': the noise power at 0.1--1~MHz (0.6), the
end-of-cycle tail level (1.0), the droop half-time (2.0), the skips 1--4 ringing (2.4), the baseline droop amplitude (2.4), the CDS autocorrelation at lag 2 (2.4) and
the correlated noise floor (2.6). The first-skip ringing frequency has no run-to-run variation at all (with the original model the fit sits at
its lowest frequency in every pixel; corrected on 2026-10-13, see the model above) and is excluded from the ratio. Judging every pair against this spread (right panel of Figure~\ref{fig:verdicts}) flags 0 of the 35 tools for each
of the six same-amplifier pairs and 22 to 28 for the cross-amplifier pairs. The same-amplifier count is optimistic because the spread is estimated
from the same runs.

\begin{figure}[htbp]
  \centering
  \includegraphics[width=\textwidth]{../plots/compare_z_by_tool.png}
  \caption{Left: $z$ of every tool for pairs of channel-0 repeats (grey) and for channel 0 against channel 3 (red); the cross-amplifier values are
  larger for most tools. Right: the size of the channel-0 to channel-3 difference in units of the run-to-run spread of channel 0
  (log scale); bars above 3 (dashed line) are detectable.}
  \label{fig:z}
\end{figure}

\begin{figure}[htbp]
  \centering
  \includegraphics[width=0.95\textwidth]{../plots/compare_verdicts.png}
  \caption{Number of tools flagged as different for each pair of traces. Left: the pixel-level test of \code{compare\_trace.py}. Right: the same
  count judged against the run-to-run spread.}
  \label{fig:verdicts}
\end{figure}

\begin{figure}[htbp]
  \centering
  \includegraphics[width=\textwidth]{../plots/compare_waveforms.png}
  \caption{A: mean cycle shape of every run (grey: channel 0, red: channel 3). B: zoom on the pulse. C: deviation of every run from the channel-0
  mean shape at every time in the cycle (log scale); the channel-3 deviations exceed the channel-0 repeats by one to two orders of magnitude at the edges.}
  \label{fig:waveforms}
\end{figure}

\paragraph{What this means for parameter scans.} A parameter change should be judged against the run-to-run spread, which needs several repeats of the
reference in every scan. With the numbers above, a change of roughly three times that spread is detectable: about 0.5\% of the pulse peak
height or the dip height. Channel-3 data show the amplifier difference is much larger than the repeat noise, so the tools are
sensitive; whether parameter changes move them by more than 0.5\% is the open question.

\section{Limitations}

\begin{itemize}
  \item \textbf{No parameter test yet.} All traces analysed so far have identical operating parameters.
        Whether any fit responds to a parameter change is unknown.
  \item \textbf{Windows tied to this timing.} The search ranges and the cut points (dip, pulse, windows) are
        found from the data but assume the present sequencer structure. They would need review if the
        sequencer changes.
  \item \textbf{Few replicates.} Ten pixels per trace limit the smallest detectable difference and the
        smallest reportable $p$.
  \item \textbf{Fits are descriptive.} They are not validated against a circuit model; a good fit does not by
        itself prove the physical interpretation given above.
  \item \textbf{Repeats.} Only four healthy channel-0 repeats are available, so the run-to-run spread itself is poorly known.
  \item \textbf{Back-to-back acquisitions.} The CDS level changed between consecutive runs; reference repeats should be spread through a scan.
  \item \textbf{Units.} Everything is in uncalibrated ADU. Conversion to electrons needs the gain, which is one of
        the quantities the project wants to predict.
\end{itemize}

\section{Where the code lives}


\begin{center}
\begin{tabular}{ll}
\toprule
Item & Location \\
\midrule
Exponential fit        & \code{fit\_exponential} in \code{src/compare\_trace.py} \\
Damped ringing fit     & \code{fit\_damped\_ringing} in \code{src/compare\_trace.py} \\
Noise-model fit        & \code{fit\_noise\_model} in \code{src/compare\_trace.py} \\
Per-pixel measurement  & \code{skip\_tools} in \code{src/compare\_trace.py} \\
Permutation test       & \code{permutation\_test}, \code{run\_tests} in \code{src/compare\_trace.py} \\
Decoding and regions   & \code{src/decode\_trace.py}, \code{src/trace\_regions.py} \\
Figures                & \code{src/plot\_tool\_figures.py} (annotated tools, repeat comparison) \\
\bottomrule
\end{tabular}
\end{center}

\noindent Usage: \code{python compare\_trace.py TRACE\_A TRACE\_B [--perms N] [--alpha a] [--csv out.csv]}.
Figures referred to above (\code{regions\_templates.png}, \code{skip\_evolution.png}) are in the
\code{plots} folder.

\appendix
\section{Complete results: ACM 101, channel 0 against channel 3}

All 39 tools, grouped by region and ordered by $p$. The threshold for ``distinguishes'' is
$p<2.9\times10^{-4}$ (35 tools, $\alpha=0.01$); ``marginal'' means $p<0.05$ but above the threshold;
vertical-block tools are information only. $z$ is the Welch statistic described in Section~\ref{sec:stats};
for the fit time constants and other very stable quantities it is inflated, so rank by $p$.

\begin{center}
\small
\begin{tabular}{llrrl}
\toprule
Region & Tool & $z$ & $p$ & Verdict \\
\midrule
skip & pulse peak height & 284.99 & $<10^{-4}$ & distinguishes \\
skip & dip height & 231.32 & $<10^{-4}$ & distinguishes \\
skip & noise power 3-7.5 MHz & 158.10 & $<10^{-4}$ & distinguishes \\
skip & noise power 1-3 MHz & 113.72 & $<10^{-4}$ & distinguishes \\
skip & mean shape & 82.52 & $<10^{-4}$ & distinguishes \\
skip & CDS mean (skips $\ge$ 10) & 15.80 & $<10^{-4}$ & distinguishes \\
skip & signal-window level & 15.56 & $<10^{-4}$ & distinguishes \\
skip & noise power 0.1-1 MHz & 12.90 & $<10^{-4}$ & distinguishes \\
skip & CDS rms (detrended) & 9.10 & $<10^{-4}$ & distinguishes \\
skip & first-skip CDS excess & 5.92 & $<10^{-4}$ & distinguishes \\
skip & noise of 10-skip average & 4.45 & $3.0\times10^{-4}$ & marginal \\
skip & skips 1-4 CDS ringing & 2.06 & 0.053 & no \\
skip & droop half-time (skips) & 1.72 & 0.12 & no \\
skip & noise of 50-skip average & 1.45 & 0.16 & no \\
skip & baseline droop curve & 1.24 & 0.27 & no \\
skip & droop total & 1.20 & 0.3 & no \\
skip & end-of-cycle tail level & 1.15 & 0.32 & no \\
skip & baseline at skip 1 & 1.15 & 0.33 & no \\
skip & CDS autocorrelation at lag 2 & 0.85 & 0.4 & no \\
fit & dip recovery amplitude & 496.74 & $<10^{-4}$ & distinguishes \\
fit & settling amplitude after pulse & 299.07 & $<10^{-4}$ & distinguishes \\
fit & white noise per skip (ADU) & 8.53 & $<10^{-4}$ & distinguishes \\
fit & dip recovery time (us) & 385.44 & $2.0\times10^{-4}$ & distinguishes \\
fit & baseline droop time constant (skips) & 5.54 & $2.0\times10^{-4}$ & distinguishes \\
fit & settling time after pulse (us) & 112.90 & $3.0\times10^{-4}$ & marginal \\
fit & first-skip ringing decay (skips) & 4.33 & $7.0\times10^{-4}$ & marginal \\
fit & first-skip ringing amplitude & 4.23 & $1.8\times10^{-3}$ & marginal \\
fit & first-skip ringing frequency (cyc/skip) & 4.20 & $4.5\times10^{-3}$ & marginal \\
fit & baseline droop amplitude & 1.12 & 0.35 & no \\
fit & correlated noise floor (ADU) & 0.39 & 0.76 & no \\
horizontal & maximum & 31.62 & $<10^{-4}$ & distinguishes \\
horizontal & mean shape & 26.75 & $<10^{-4}$ & distinguishes \\
horizontal & end level & 14.52 & $<10^{-4}$ & distinguishes \\
horizontal & minimum & 13.38 & $<10^{-4}$ & distinguishes \\
horizontal & scatter between transfers & 0.52 & 0.58 & no \\
vertical & mean shape & 17.69 & $<10^{-4}$ & info only \\
vertical & maximum & 10.23 & $<10^{-4}$ & info only \\
vertical & minimum & 3.41 & $9.0\times10^{-4}$ & info only \\
vertical & end level & 0.05 & 0.96 & info only \\
\bottomrule
\end{tabular}
\end{center}

\noindent Source: \code{plots/compare\_ch0\_vs\_ch3\_acm101.csv}, written by
\code{python compare\_trace.py TRACE\_A TRACE\_B --csv ...}. Both split controls on trace A passed.

\end{document}
